\section{Related Work}

\paragraph{Neural codec language models.}
Discrete neural audio codecs such as SoundStream
\cite{zeghidour2021soundstream} and EnCodec \cite{defossez2022encodec}
established low-bitrate residual vector-quantized speech representations, and
VALL-E demonstrated that a language model over such codes can perform
text-to-speech \cite{wang2023valle}. Later systems emphasize streaming
generation from codec tokens, including Moshi's full-duplex speech modeling
\cite{defossez2024moshi} and CosyVoice~2's streaming synthesis
\cite{du2024cosyvoice2}. Description-conditioned control of voice attributes
through natural-language annotations was explored by Lyth and King
\cite{lyth2024parler}; Qwen3-TTS VoiceDesign exposes the same control paradigm.
These works motivate the model family; this paper contributes only the serving
implementation.

\paragraph{Qwen3-TTS.}
The Qwen3-TTS technical report presents a family of multilingual,
description-controllable text-to-speech models and two speech-tokenizer rates
\cite{hu2026qwen3tts}. The 12-Hz design uses residual vector-quantized speech
tokens: a backbone predicts the zeroth codebook and a multi-token prediction
module produces the residual codebooks. The report evaluates model quality and
streaming efficiency under its own inference environment. Our work neither
changes the trained model nor reproduces those model-quality claims. It studies
a separate systems question: how to execute the released 1.7B VoiceDesign
checkpoint through a purpose-built Rust/CUDA path on DGX Spark.

\paragraph{LLM serving systems and the stock comparator.}
General-purpose LLM serving frameworks optimize throughput with techniques
such as paged key-value memory in vLLM \cite{kwon2023vllm} and RadixAttention
in SGLang \cite{zheng2023sglang}. These systems target broad model coverage
inside a Python runtime; \projectname{} instead specializes one checkpoint and
one GPU architecture in a framework-free native binary. The comparison in this paper uses the separately
versioned SGLang-Omni 0.1.0 Qwen3-TTS serving path, not the experiments reported
in the original SGLang paper. Its upstream source commit, container base,
dependency inventory, compatibility patch, and launch configuration are pinned
in the benchmark evidence. The comparator remains labeled \emph{stock} only
when packaging changes do not alter scheduling, model execution, codec cadence,
or streaming behavior.

\paragraph{Scope of comparison.}
The paper compares serving implementations, not trained model quality. Both
subjects use the same checkpoint repository, immutable revision, ordered input
workload, language values, VoiceDesign instructions, seeds, portable sampling
controls, duration ceiling, client, host, and audio measurement boundary.
Implementation-specific model artifacts are recorded separately because the
native runtime packages only decoder inference material whereas the stock stack
may retain a broader speech-tokenizer artifact.
